\begin{textsample}{POS Dim 2 – vem\_ed\_19 – Score 15.00 – vem\_ed\_19/t090.txt}  \label{ex:f2_pos_008}
continuação Reflexões: gestão dos projetos Em 13 de \textbf{junho}, Ana Carolina Freschi, do apoio aos PCIs em inglês na equipe PCI / Cesu, apresentou o trabalho com seus colegas: Divinia Jithoo, Lindelwa Mkhize (Durban University of Technology, África do Sul), Daniel Otieno Okech (Kenyatta University, Quênia), Eva Haug (AUAS, Holanda) e Harshita Tripati (Shiv Nadar University, Índia).
Em pauta, a \textbf{dinâmica} de facilitação de Intercâmbios Virtuais nos diversos contextos educacionais desses países — sob os pontos de vista tecnológico, intercultural, acadêmico e linguístico.
Nesse mesmo dia, o painel “¡ SOS!
Intervención del Coordinador COIL” foi \textbf{conduzido} por Succi Junior, Gisselle Morales Veloquio (diretora de aprendizagem global do Tecnológico de Monterrey, México) e Regiane Moreira (professora responsável pelo apoio aos PCIs em espanhol nas Fatecs).
Os autores sugerem 10 \textbf{passos} para a \textbf{intervenção} do coordenador de Intercâmbios Virtuais, organizados nas \textbf{seguintes} etapas: \textbf{institucionalizar}; \textbf{definir} claramente o papel do coordenador COIL; alinhar expectativas da instituição, seus coordenadores, professores e parceiros; determinar \textbf{critérios} para \textbf{busca} de parceiros; oferecer formação continuada (para professores novos e experientes em Intercâmbios Virtuais); \textbf{gerenciar} o \textbf{design} dos projetos; realizar o acompanhamento dos projetos; criar índices de \textbf{comparação} e de informação; \textbf{explicar} e difundir constantemente os projetos COIL; decidir quando intervir (\textbf{intervenção} de \textbf{emergência}, apenas quando necessário).

% matched lemmas: busca, comparação, conduzir, critério, definir, design, dinâmica, emergência, explicar, gerenciar, institucionalizar, intervenção, junho, passo, seguinte
\end{textsample}
